\section{Discussion and conclusion}

We set out to answer a simple-sounding question --- how much of the Sun's
energy can humanity take --- and found that the honest answer inverts the usual
framing. The resource is astronomically abundant (\InterceptedTW~TW), land is
not scarce (world electricity fits on a \AreaWorldSide~km square of desert
panel), and even orbit, the ``layer above the atmosphere,'' is a
capacity-factor improvement rather than a new source of free power. The one
real, physical constraint is that the light reaching the ground is already
running the biosphere.

The contribution we are most confident in is turning that constraint into an
optimisation. Because an ecosystem's light requirement is a daily
\emph{integral}, harvesting has temporal slack, and exploiting it --- yielding
light to life at low-value hours and harvesting the peaks --- recovers
\TemporalGainHead\% more usable energy at $f_{\min}=\FminHead$ and up to
\TemporalGainMax\% at higher reserves, all at zero additional ecological cost.
The optimum is a transparent water-filling rule, not a black box.

\paragraph{Limitations.} The model is deliberately coarse and self-contained:
latitude bands stand in for biomes, the ecosystem constraint is a single
daily-integral floor rather than a species- and process-specific light
response, and demand is a stylised duck curve. The \texttt{gpu\_full} config
refines the spatial and temporal grids and the reserve sweep; substituting
gridded insolation and land-cover data, per-biome light-response curves, and
real regional demand would sharpen the numbers without changing the structure.
The qualitative claim --- that solar harvesting is an allocation problem
against the biosphere, and that the allocation is schedulable --- is what the
model is built to show.

\paragraph{Outlook.} The framing suggests a research programme: replace the
integral floor with measured net-primary-production light responses, couple the
harvest schedule to storage and grid dynamics, and let space-based collectors
relax the temporal constraint entirely for the hours the biosphere claims. The
Sun is not the limit; the biosphere is, and it is one we can plan around.
